\chapter*{ABSTRAK}
\addcontentsline{toc}{chapter}{ABSTRAK}

\begin{center}
    \textbf{Pengembangan GraphRAG dengan Pembobotan Graf Pengetahuan Berbasis Penilaian Pengguna untuk Mengurangi Halusinasi pada LLM} \\
    \vspace{0.4cm}
    oleh \\
    \vspace{0.4cm}
    Dartika Anie Marian\\
    23524319\\
\end{center}

\begin{spacing}{1.0}
\emph{Large Language Model} (LLM) rentan menghasilkan halusinasi, yaitu keluaran yang tampak meyakinkan tetapi tidak sesuai fakta. \emph{Retrieval-Augmented Generation} (RAG) dan \emph{Graph-based RAG} (GraphRAG) mengurangi masalah ini tanpa mengubah parameter model, tetapi keandalan sumber pengetahuan yang diambil umumnya belum dipertimbangkan. Pada pendekatan GraphRAG yang ditelaah, bobot keandalan diturunkan dari asal-usul atau otoritas dokumen, dan belum ditemukan penggunaan penilaian pengguna terhadap suatu pengetahuan sebagai dasar bobot dalam pengambilan konteks. Penelitian ini merancang mekanisme pembobotan graf pengetahuan berbasis penilaian pengguna, yaitu status \emph{accepted answer} dan skor suara, pada GraphRAG menggunakan \emph{Design Science Research Methodology}. Artefaknya adalah graf pengetahuan berbobot dan mekanisme pengambilan konteks tiga tahap yang memakai bobot tersebut. Pengujian dilakukan pada satu studi kasus, yaitu data Stack Overflow, dengan 384 pertanyaan uji, model generatif GPT-4o-mini, dan penilaian halusinasi oleh panel tiga LLM \emph{judge} tanpa validasi anotator manusia. Lima konfigurasi dibandingkan, yaitu LLM murni (A), RAG vektor (B-Vec), RAG vektor dengan bobot \emph{accepted} tanpa graf (B-Acc), GraphRAG dengan bobot \emph{accepted} (C-Acc), dan GraphRAG dengan bobot \emph{accepted} dan skor suara (C-Acc+Vote). Dengan penolakan menjawab dikeluarkan dari penyebut, \emph{hallucination rate} masing-masing adalah 25,8\%; 24,8\%; 19,4\%; 18,8\%; dan 17,0\%. Pembobotan \emph{accepted answer} pada B-Acc, tanpa graf, sudah menurunkan angka 5,4 poin terhadap B-Vec, sedangkan tambahan dari mekanisme berbasis graf hanya 0,6 poin sehingga manfaat graf belum terbukti. Pengaruh skor suara bersifat eksploratif dan belum dapat disimpulkan karena arahnya tidak konsisten antar-\emph{judge}. Konteks berbobot memuat lebih banyak jawaban \emph{accepted}, tetapi relevansi konteks dan \emph{faithfulness} sedikit lebih rendah daripada B-Vec. Terhadap LLM murni, penurunan berbalik bila penolakan menjawab dihitung sebagai kegagalan. Kontribusi penelitian ini adalah (i) mekanisme pembobotan graf pengetahuan berbasis penilaian pengguna beserta implementasinya, (ii) protokol evaluasi yang memisahkan pengaruh pembobotan dari struktur graf, melaporkan penolakan menjawab, dan mencegah kebocoran data uji, serta (iii) bukti empiris bahwa bobot \emph{accepted answer} berperan dalam menurunkan halusinasi. Hasil berlaku pada satu studi kasus, satu \emph{run} per konfigurasi, dan label panel \emph{judge} yang sangat berbeda tingkat kelonggarannya, sehingga selisih bersifat deskriptif dan tidak digeneralisasi.

\vspace{0.3cm}
\noindent Kata kunci: halusinasi LLM, GraphRAG, penilaian pengguna, pembobotan graf pengetahuan, \emph{retrieval-augmented generation}, Stack Overflow
\end{spacing}